\documentclass[12pt,]{article}
\usepackage{lmodern}
\usepackage{amssymb,amsmath}
\usepackage{ifxetex,ifluatex}
\usepackage{fixltx2e} % provides \textsubscript
\ifnum 0\ifxetex 1\fi\ifluatex 1\fi=0 % if pdftex
  \usepackage[T1]{fontenc}
  \usepackage[utf8]{inputenc}
\else % if luatex or xelatex
  \ifxetex
    \usepackage{mathspec}
  \else
    \usepackage{fontspec}
  \fi
  \defaultfontfeatures{Ligatures=TeX,Scale=MatchLowercase}
\fi
% use upquote if available, for straight quotes in verbatim environments
\IfFileExists{upquote.sty}{\usepackage{upquote}}{}
% use microtype if available
\IfFileExists{microtype.sty}{%
\usepackage{microtype}
\UseMicrotypeSet[protrusion]{basicmath} % disable protrusion for tt fonts
}{}
\usepackage[margin=1in]{geometry}
\usepackage{hyperref}
\hypersetup{unicode=true,
            pdftitle={Fitting Distances by Least Squares},
            pdfauthor={Jan de Leeuw},
            pdfborder={0 0 0},
            breaklinks=true}
\urlstyle{same}  % don't use monospace font for urls
\usepackage{graphicx,grffile}
\makeatletter
\def\maxwidth{\ifdim\Gin@nat@width>\linewidth\linewidth\else\Gin@nat@width\fi}
\def\maxheight{\ifdim\Gin@nat@height>\textheight\textheight\else\Gin@nat@height\fi}
\makeatother
% Scale images if necessary, so that they will not overflow the page
% margins by default, and it is still possible to overwrite the defaults
% using explicit options in \includegraphics[width, height, ...]{}
\setkeys{Gin}{width=\maxwidth,height=\maxheight,keepaspectratio}
\IfFileExists{parskip.sty}{%
\usepackage{parskip}
}{% else
\setlength{\parindent}{0pt}
\setlength{\parskip}{6pt plus 2pt minus 1pt}
}
\setlength{\emergencystretch}{3em}  % prevent overfull lines
\providecommand{\tightlist}{%
  \setlength{\itemsep}{0pt}\setlength{\parskip}{0pt}}
\setcounter{secnumdepth}{5}
% Redefines (sub)paragraphs to behave more like sections
\ifx\paragraph\undefined\else
\let\oldparagraph\paragraph
\renewcommand{\paragraph}[1]{\oldparagraph{#1}\mbox{}}
\fi
\ifx\subparagraph\undefined\else
\let\oldsubparagraph\subparagraph
\renewcommand{\subparagraph}[1]{\oldsubparagraph{#1}\mbox{}}
\fi

%%% Use protect on footnotes to avoid problems with footnotes in titles
\let\rmarkdownfootnote\footnote%
\def\footnote{\protect\rmarkdownfootnote}

%%% Change title format to be more compact
\usepackage{titling}

% Create subtitle command for use in maketitle
\providecommand{\subtitle}[1]{
  \posttitle{
    \begin{center}\large#1\end{center}
    }
}

\setlength{\droptitle}{-2em}

  \title{Fitting Distances by Least Squares}
    \pretitle{\vspace{\droptitle}\centering\huge}
  \posttitle{\par}
    \author{Jan de Leeuw}
    \preauthor{\centering\large\emph}
  \postauthor{\par}
      \predate{\centering\large\emph}
  \postdate{\par}
    \date{First created on May 21, 2019. Last update on August 21, 2019}


\begin{document}
\maketitle
\begin{abstract}
We review the continuity and differentiability properties of the stress
loss function in MDS and the local and global convergence properties of
the SMACOF algorithm.
\end{abstract}

{
\setcounter{tocdepth}{3}
\tableofcontents
}
\textbf{Note:} This is a working paper which will be expanded/updated
frequently. All suggestions for improvement are welcome. The directory
\href{http://deleeuwpdx.net/pubfolders/statioclass}{deleeuwpdx.net/pubfolders/statioclass}
has a pdf version, the bib file, the complete Rmd file with the code
chunks, and the R source code.

\hypertarget{introduction}{%
\section{Introduction}\label{introduction}}

This book is about the least squares loss function
\begin{equation}\label{E:stress}\sigma(X)=\frac12\mathop{\sum}_{1\leq i<j\leq n} w_{ij}(\delta_{ij}-d_{ij}(X))^2.
\end{equation} define on \(\mathbb{R}^{n\times p}\), the space of all
\(n\times p\) matrices. We follow Kruskal (1964) and call \(\sigma(X)\)
the \emph{stress} of \emph{configuration} \(X\). Minimizing stress over
p-dimensional configurations is the \emph{pMDS problem}.

In \(\eqref{E:stress}\) the matrices of \emph{weights} \(W=\{w_{ij}\}\)
and \emph{dissimilarities} \(\Delta=\{\delta_{ij}\}\) are symmetric,
non-negative, and hollow (zero diagonal). The matrix-valued function
\(D(X)=\{d_{ij}(X)\}\) contains Euclidean distances between the rows of
the configuration \(X\), which are the coordinates of \(n\)
\emph{points} in \(\mathbb{R}^p\). Thus \(D(X)\) is also symmetric,
non-negative, and hollow.

Two important special cases of pMDS are \emph{Unidimensional Scaling},
which is 1MDS, and \emph{Full-dimensional Scaling}, which is nMDS.
Because of their importance, they also have their very own acronyms UDS
and FDS, and they have their own chapters in this book.

In pMDS we minimize stress. This means that we are trying to find the
\emph{global minimum} over p-dimensional configurations. In practice,
however, our algorithms find \emph{local minima}, which may or may not
be global. In this paper we will summarize what we know about global and
local minima, and what we know about local maxima and saddle points.

\hypertarget{notation}{%
\section{Notation}\label{notation}}

First some convenient notation, first introduced in De Leeuw (1977).
Vector \(e_i\) has \(n\) elements, with element \(i\) equal to \(+1\),
and all other elements zero. \(A_{ij}\) is the matrix
\((e_i-e_j)(e_i-e_j)'\), which means elements \((i,i)\) and \((j,j)\)
are equal to \(+1\), while \((i,j)\) and \((j,i)\) are \(-1\). Thus
\begin{equation}\label{E:anot}
d_{ij}^2(X)=(e_i-e_i)'XX'(e_i-e_j)=\text{tr}\ X'A_{ij}X=\text{tr}\ A_{ij}C,
\end{equation} with \(C=XX'\).

We also define \begin{equation}\label{E:V}
V=\mathop{\sum}_{1\leq i<j\leq n}w_{ij}A_{ij},
\end{equation} and the matrix-valued function \(B(\bullet)\) with
\begin{equation}\label{E:B}
B(X)=\mathop{\sum}_{d_{ij}(X)>0}w_{ij}\frac{\delta_{ij}}{d_{ij}(X)}A_{ij},
\end{equation} and \(B(X)=0\) if \(X=0\). If we assume, without loss of
generality, that \[
\frac12\mathop{\sum}_{1\leq i<j\leq n}w_{ij}\delta_{ij}^2=1,
\] then \begin{equation}\label{E:S}
\sigma(X)=1-\text{tr}\ X'B(X)X+\frac12\text{tr}\ X'VX.
\end{equation}

We also suppose, without loss of generality, that \(W\) is
\emph{irreducible}, so that the pMDS problem does not separate into a
number of smaller pMDS problems. For symmetric matrices irreducibity
means that we cannot find a permutation matrix \(\Pi\) such that
\(\Pi'W\Pi\) is the direct sum of a number of smaller matrices.

\(V\) is symmetric with non-positive off-diagonal elements. It is
doubly-centered (rows and columns add up to zero) and thus weakly
diagonally dominant. It follows that it is positive semi-definite (see
Varga (1962), section 1.5). Because of irreducibility it has rank
\(n-1\), and the vectors in its null space are all proportional to
\(e\), the vector with all elements equal to \(+1\). The matrix \(B(X)\)
is also symmetric, positive semi-definite, and doubly-centered for each
\(X\). It may not be irreducible, because for example \(B(0)=0\).

\hypertarget{differentiability}{%
\section{Differentiability}\label{differentiability}}

Obviously \[
\mathcal{D}\eta^2(X)=2VX
\] Additionally \(\rho(\bullet)\) is differentiable at \(X\) if and only
if \(d_{ij}(X)>0\) for all \(i<j\) with \(w_{ij}\delta_{ij}>0\). In that
case \[
\mathcal{D}\rho(X)=B(X)X.
\] \textbf{Result 13: } \(\sigma(\bullet)\) is differentiable at at
\(X\) if and only if \(d_{ij}(X)>0\) for all \(i<j\) with
\(w_{ij}\delta_{ij}>0\). In that case \[
\mathcal{D}\sigma(X)=(V-B(X))X.
\] \textbf{Result 14: } If \(\sigma(\bullet)\) is differentiable at the
stationary point \(X\) then \((V-B(X))X=0\), or, in fixed point form,
\(X=V^+B(X)X\).

\hypertarget{dc-functions}{%
\section{DC Functions}\label{dc-functions}}

Following De Leeuw (1977) we also define \begin{align}
\rho(X)&=\mathop{\sum}_{1\leq i<j\leq n} w_{ij}\delta_{ij}d_{ij}(X)=\text{tr}\ X'B(X)X,\\
\eta^2(X)&=\mathop{\sum}_{1\leq i<j\leq n} w_{ij}d_{ij}^2(X)=\text{tr}\ X'VX.
\end{align}

\textbf{Result 1: {[}DC{]}}

\begin{enumerate}
\def\labelenumi{\arabic{enumi}.}
\tightlist
\item
  \(\rho(\bullet)\) is convex, homogeneous of degree one, non-negative,
  and continuous.
\item
  \(\eta^2(\bullet)\) is convex, homogeneous of degree two,
  non-negative, quadratic, and continuous.
\item
  \(\sigma(\bullet)\) is a difference of convex functions (a.k.a. a
  \emph{DC function} or a \emph{delta-convex} function).
\end{enumerate}

See Hirriart-Urruty (1988) or Vesely and Zajicek (1989) for a general
discussion of DC functions.

From the literature we know \(\sigma(\bullet)\) is locally Lipschitz,
has directional derivatives and nonempty subdifferentials everywhere,
and is twice differentiable almost everywhere. In fact
\(\sigma(\bullet)\) is infinitely many times differentiable at all \(X\)
that have \(d_{ij}(X)>0\) for all \(i<j\) with \(w_{ij}\delta_{ij}>0\).

\hypertarget{homogeneity}{%
\section{Homogeneity}\label{homogeneity}}

\textbf{Result 2: {[}Bounded{]}} At a local minimum point \(X\) of
\(\sigma(\bullet)\) we have \(\eta(X)\leq 1\).

\textbf{Proof:} Homogeneity gives
\(\sigma(\lambda X)=1-\lambda\rho(X)+\frac12\lambda^2\eta^2(X)\). If
\(X\) is a local minimum then the minimum over \(\lambda\) is attained
for \(\lambda=1\), i.e.~we must have \(\rho(X)=\eta^2(X)\). By
Cauchy-Schwartz \(\rho(X)\leq\eta(X)\) for all \(X\), and thus at a
local minimum \(\eta^2(X)\leq\eta(X)\), i.e. \(\eta(X)\leq 1\).
\(\blacksquare\)

\textbf{Result 3: {[}Local Maxima{]}} \(X=0\) is the unique local
maximum point of \(\sigma(\bullet)\) and \(\sigma(0)=1\) is the unique
local maximum.

\textbf{Proof:} This is because
\(\sigma(\lambda X)=1-\lambda\rho(X)+\frac12\lambda^2\eta^2(X)\) is a
convex quadratic in \(\lambda\). The only maximum on the ray through
\(X\) occurs at the boundary \(\lambda=0\). \(\blacksquare\)

\textbf{Result 4: {[}Unbounded{]}} \(\sigma(\bullet)\) is unbounded and
consequently has no global maximum.

\textbf{Proof:} \(\blacksquare\)

Note that the unboundedness of stress also follows from the proof of
result 3.

\hypertarget{picture}{%
\section{Picture}\label{picture}}

Here is a picture of stress on a two-dimensional subspace which
illustrates both result 3 and result 15.

We first make a global perspective plot, over the range \((-2.5,+2.5)\).

\begin{center}\includegraphics{statioclass_files/figure-latex/global_perspective-1} \end{center}

Figure 1: Picture of Stress

Note the sharp ridge going northwest-southeast in the plot, indicating a
ray of configurations where one of more of the distances are zero, and
where consequently \(\sigma(\bullet)\) is not differentiable.

We first make a global perspective plot, over the range \((-2.5,+2.5)\).

\begin{center}\includegraphics{statioclass_files/figure-latex/global_perspective2-1} \end{center}

Figure 1: Picture of Stress

We first make a global perspective plot, over the range \((-2.5,+2.5)\).

\begin{center}\includegraphics{statioclass_files/figure-latex/global_perspective3-1} \end{center}

Figure 1: Picture of Stress

\hypertarget{one-sided-directional-derivatives}{%
\section{One-sided Directional
Derivatives}\label{one-sided-directional-derivatives}}

For a function \(f:\mathbb{R}^{n\times p}\rightarrow\mathbb{R}\) we
define the one-sided (Dini) directional derivative at \(X\) in direction
\(Y\) as \[
df(X;Y)=\lim_{\epsilon\downarrow 0}\frac{f(X+\epsilon Y)-f(X)}{\epsilon},
\] and the one-sided (Peano) second-order directional derivative at
\(X\) in direction \(Y\) as \[
d^2f(X;Y)=\lim_{\epsilon\downarrow 0}\frac{f(X+\epsilon Y)-f(X)-\epsilon\ \mathcal{D}f(X;Y)}{\frac12\epsilon^2}.
\] We know that \(\sigma(\bullet)\) has one-sided directional
derivatives of orders one and two. These can be easily computed by
expanding the function at \(X\) in direction \(Y\).

\hypertarget{expansion}{%
\subsection{Expansion}\label{expansion}}

We start by expanding squared distances and distances. First
\begin{equation}\label{E:dd}
d_{ij}^2(X+\epsilon Y)=d_{ij}^2(X)+2\epsilon\ \text{tr}\ Y'A_{ij}X+\epsilon^2\ \text{tr}\ Y'A_{ij}Y.
\end{equation} Then for \(d_{ij}(X)>0\) \begin{equation}\label{E:dg}
d_{ij}(X+\epsilon Y)=d_{ij}(X)+\epsilon\frac{\text{tr}\ Y'A_{ij}X}{d_{ij}(X)}+
\frac12\epsilon^2\frac{1}{d_{ij}(X)}\left\{d_{ij}^2(Y)
-\frac{(\text{tr}\ Y'A_{ij}X)^2}{\text{tr}\ X'A_{ij}X}\right\}+o(\epsilon^2),
\end{equation} and for \(d_{ij}(X)=0\) \begin{equation}\label{E:d0}
d_{ij}(X+\epsilon Y)=\epsilon\ d_{ij}(Y).
\end{equation} Equations \(\eqref{E:dd}\), \(\eqref{E:dg}\), and
\(\eqref{E:d0}\) combine in a straightforward way to an expansion of
\(\sigma(\bullet)\) at \(X\) in direction \(Y\).

\textbf{Result 5: {[}Quadratic{]}} \begin{multline}
\sigma(X+\epsilon Y)=\sigma(X)+\epsilon\left\{\text{tr}\ Y'(V-B(X))X-\sum_{d_{ij}(X)=0}w_{ij}\delta_{ij}d_{ij}(Y)\right\}\\ 
+\frac12\epsilon^2\left\{\text{tr}\ Y'(V-B(X))Y+\sum_{d_{ij}(X)>0}w_{ij}\frac{\delta_{ij}}{d_{ij}(X)}
\frac{(\text{tr}\ Y'A_{ij}X)^2}{\text{tr}\ X'A_{ij}X}\right\}+o(\epsilon^2)
\end{multline}

There are some interesting special cases of this result.

\textbf{Result 6: {[}Antisymmetric{]}} Suppose \(Y=XT\), with \(T\)
antisymmetric, so that \(X+\epsilon Y=X(I+\epsilon T)\). Then
\begin{equation}
\sigma(X+\epsilon Y)=\sigma(X)-\epsilon\sum_{d_{ij}(X)=0}w_{ij}\delta_{ij}d_{ij}(Y)
+\frac12\epsilon^2\text{tr}\ Y'(V-B(X))Y+o(\epsilon^2)
\end{equation}

\textbf{Result 7: {[}Singular{]}} Suppose \(\underline{X}=[X\mid 0]\)
and \(\underline{X}=[0\mid Y]\) so that
\(\underline{X}+\epsilon\underline{Y}=[X\mid\epsilon Y]\). Here
\(\underline{X}\) and \(\underline{Y}\) are \(n\times p\), \(X\) is
\(n\times r\), with \(r<p\), and \(Y\) is \(n\times(p-r)\). Then
\begin{equation}
\sigma(\underline{X}+\epsilon \underline{Y})=\sigma(X)-\epsilon\sum_{d_{ij}(X)=0}w_{ij}\delta_{ij}d_{ij}(Y)
+\frac12\epsilon^2\text{tr}\ Y'(V-B(X))Y+o(\epsilon^2)
\end{equation}

\textbf{Result 7: {[}Singular{]}} Suppose \(\underline{X}=[X\mid 0]\)
and \(\underline{X}=[Z\mid Y]\) so that
\(\underline{X}+\epsilon\underline{Y}=[X+\epsilon Z\mid\epsilon Y]\).
Here \(\underline{X}\) and \(\underline{Y}\) are \(n\times p\), \(X\) is
\(n\times r\), with \(r<p\), and \(Y\) is \(n\times(p-r)\). Then
\begin{equation}
\sigma(\underline{X}+\epsilon \underline{Y})=\sigma(X)-\epsilon\sum_{d_{ij}(X)=0}w_{ij}\delta_{ij}d_{ij}(Y)
+\frac12\epsilon^2\text{tr}\ Y'(V-B(X))Y+o(\epsilon^2)
\end{equation}

\hypertarget{derivatives}{%
\subsection{Derivatives}\label{derivatives}}

The expansion in result 5 immediately gives the first-order and
second-order directional derivatives.

\textbf{Result 8: {[}First Directional Derivatives{]}} \begin{equation}
d\sigma(X;Y)=
\text{tr}\ X'(V-B(X))Y-\sum_{d_{ij}(X)=0}w_{ij}\delta_{ij}d_{ij}(Y).
\end{equation}

\textbf{Result 9: {[}Second Directional Derivatives{]}} \begin{equation}
d^2\sigma(X;Y)=
\text{tr}\ Y'(V-B(X))Y+\sum_{d_{ij}(X)>0}w_{ij}\frac{\delta_{ij}}{d_{ij}(X)}\frac{(\text{tr}\ Y'A_{ij}X)^2}{\text{tr}\ X'A_{ij}X}.
\end{equation}

The expression for the second order derivatives can be made a bit more
matrix and computer friendly by rewriting it differently. Define
\(y=\text{vec}(Y)\) and \(x=\text{vec}(X)\). Thus \(x\) and \(y\) have
length \(np\). Also define
\(\overrightarrow{A}_{ij}=I_p\otimes A_{ij}\), with \(I_p\) the identity
matrix of order \(p\) and \(\otimes\) the Kronecker product. Thus the
\(np\times np\) matrix \(\overrightarrow{A}_{ij}\) is the direct sum of
\(p\) copies of our previous \(n\times n\) matrix \(A_{ij}\). In the
same way we write \(\overrightarrow{V}\) for \(I_p\otimes V\) and
\(\overrightarrow{B}(X)\) for \(I_p\otimes B(X)\).

With this new notation \begin{equation}
d^2\sigma(X;Y)=y'\left\{\overrightarrow{V}-\overrightarrow{B}(X)+\mathop{\sum}_{d_{ij}(X)>0}w_{ij}\frac{\delta_{ij}}{d_{ij}(X)}\frac{\overrightarrow{A}_{ij}xx'\overrightarrow{A}_{ij}}{x'\overrightarrow{A}_{ij}x}\right\}y.
\end{equation} Define, for some additional shorthand, the
\(np\times np\) matrix \[
G(X)=\sum_{d_{ij}(X)>0}w_{ij}\frac{\delta_{ij}}{d_{ij}(X)}\frac{\overrightarrow{A}_{ij}xx'\overrightarrow{A}_{ij}}{x'\overrightarrow{A}_{ij}x},
\] as well as \(H(X)=\overrightarrow{B}(X)-G(X)\). Note that \[
H(X)=\sum_{d_{ij}(X)>0}w_{ij}\frac{\delta_{ij}}{d_{ij}(X)}\left\{\overrightarrow{A}_{ij}-\frac{\overrightarrow{A}_{ij}xx'\overrightarrow{A}_{ij}}{x'\overrightarrow{A}_{ij}x}\right\},
\] and consequently both \(G(X)\) and \(H(X)\) are positive
semi-definite, with \(H(X)\) singular because \(H(X)x=0\). This implies
\[
\text{tr}\ Y'(V-B(X))Y\leq d^2\sigma(X;Y)\leq\text{tr}\ Y'VY.
\] Also note that if \(d_{ij}(X)>0\) for all \(i<j\) then
\(\sigma(\bullet)\) is two times (Fréchet) differentiable at \(X\), with
\begin{equation}
\mathcal{D}\sigma(X)=(V-B(X))X.
\end{equation} We have to be a bit careful with the formula for the
second (Fréchet) derivative, which is a map for which there is no
straightforward matrix expression. Its value at \(Y\) is
\begin{equation}
\mathcal{D}^2\sigma(X)(Y)=d^2\sigma(X;Y)=y'(\overrightarrow{V}-H(X))y,
\end{equation} with \(y=\text{vec}(Y)\) as usual.

Note that result 6 imply that

\hypertarget{necessary-conditions}{%
\section{Necessary Conditions}\label{necessary-conditions}}

The formulas for the first and second order one-sided directional
derivatives can be used to give necessary conditions for a local minimum
(see, for example, Bednarik and Pastor (2008) or Ivanov (2016)).

\textbf{Result 10: {[}Local Minimum 1{]}} If \(X\) is a local minimum
point of \(\sigma(\bullet)\) then

\begin{enumerate}
\def\labelenumi{\arabic{enumi}.}
\tightlist
\item
  \(d_{ij}(X)>0\) for all \(i<j\) with \(w_{ij}\delta_{ij}>0\),
\item
  \((V-B(X))X = 0\).
\end{enumerate}

\textbf{Proof:} This follows from 8. Suppose \((V-B(X))X\not= 0\). Then
we can find \(Y\) such that \(\text{tr}\ Y'(V-B(X))X<0\), and
\(d\sigma(X;Y)<0\). Thus \(X\) is not a local minimum point. If
\((V-B(X))X=0\) and there is an \(i<j\) such that
\(w_{ij}\delta_{ij}>0\) and \(d_{ij}(X)=0\) then again we can find \(Y\)
such that \(d\sigma(X;Y)<0\), and thus \(X\)is not a local minimum
point. \(\blacksquare\).

This result was proved for the first time in De Leeuw (1984). Note that
it implies that if \(w_{ij}\delta_{ij}>0\) for all \(i<j\) then
\(\sigma(\bullet)\) is two times (Frechet) differentiable in a
neighborhood of each local minimum \(X\), because all \(d_{ij}(X)\) must
be non-zero. But note there are situations where the result does not
apply. In multidimensional unfolding there are row-points and
column-points. Weights between two row-points and between two-column
points are zero, and thus at local minima distances between row points
and between column points can be zero, and stress is not differentiable.
Also in the case of perfect fit, if \(\delta_{ij}=d_{ij}(Z)\) for some
\(Z\), and one of more of the \(d_{ij}(Z)\) are zero, we have
\(\sigma(Z)=0\) and thus some of the distances for the global minimum,
which is clearly a local minimum, are zero.

\textbf{Result 11: {[}Local Minimum 2{]}} If \(X\) is a local minimum
point of \(\sigma(\bullet)\) then \(\overrightarrow{V}-H(X)\) is
positive semi-definite.

\textbf{Result 12: {[}Local Minimum 3{]}} If \([X\mid 0]\) is a local
minimum point of \(\sigma(\bullet)\) then

\begin{enumerate}
\def\labelenumi{\arabic{enumi}.}
\tightlist
\item
  \(d_{ij}(X)>0\) for all \(i<j\) with \(w_{ij}\delta_{ij}>0\).
\item
  \((V-B(X))X=0\).
\item
  \(V-B(X)\) is positive semi-definite.
\end{enumerate}

\textbf{Result 16: } If \(X\) is a stationary point of
\(\sigma(\bullet)\) then \(\eta(X)\leq 1\).

\textbf{Proof:} At a local minimum \((V-B(X))X=0\) which implies
\(\rho(X)=\eta^2(X)\). By Cauchy-Schwartz \(\rho(X)\leq\eta(X)\) and
thus \(\eta^2(X)\leq\eta(X)\), which implies \(\eta(X)\leq 1\).
\(\blacksquare\)

\hypertarget{full-dimensional-scaling}{%
\section{Full-dimensional Scaling}\label{full-dimensional-scaling}}

\hypertarget{unidimensional-scaling}{%
\section{Unidimensional Scaling}\label{unidimensional-scaling}}

\hypertarget{subdifferentials}{%
\section{Subdifferentials}\label{subdifferentials}}

The subdifferential of a convex function \(f(\bullet)\) at a point \(x\)
is the set \(\partial f(x)\) defined by \[
\partial f(x)=\{y\mid f(z)\geq f(x)+y'(z-x)\ \ \ \forall z\}
\] If \(f\) is differentiable at \(x\) the subdifferential is a
singleton, and \(\partial f(x)=\{\mathcal{D}f(x)\}\) (Rockafellar
(1970), Section 23). In general \[
df(x;y)=\sup\{z'y\mid  z\in\partial f(x)\}
\] \[
\partial d_{ij}^2(X)=\{2\ A_{ij}X\},
\] To compute the subdifferential \(\partial d_{ij}(X)\) we use the
representation \[
d_{ij}(X)=\max_{u'u=1}\ (e_i-e_j)'Xu.
\] From this we have for \(d_{ij}(X)>0\) \[
\partial d_{ij}(X)=\{\frac{1}{d_{ij}(X)}A_{ij}X\},
\] while for \(d_{ij}(X)=0\) \[
\partial d_{ij}(X)=\mathbb{co}\{(e_i-e_j)u'\mid u'u=1\},
\] where \(\mathbb{co}(\bullet)\) is the closed convex hull.

A stationary point of a locally Lipschitz function \(f\) is a point
\(x\) where \(0\in\partial f(x)\).

\hypertarget{smacof}{%
\section{SMACOF}\label{smacof}}

\hypertarget{references}{%
\section*{References}\label{references}}
\addcontentsline{toc}{section}{References}

\hypertarget{refs}{}
\leavevmode\hypertarget{ref-bednarik_pastor_08}{}%
Bednarik, D., and K. Pastor. 2008. ``On Second-order Conditions in
Unconstrained Optimization.'' \emph{Mathematical Programming, Series A}
113: 283--98.

\leavevmode\hypertarget{ref-deleeuw_C_77}{}%
De Leeuw, J. 1977. ``Applications of Convex Analysis to Multidimensional
Scaling.'' In \emph{Recent Developments in Statistics}, edited by J.R.
Barra, F. Brodeau, G. Romier, and B. Van Cutsem, 133--45. Amsterdam, The
Netherlands: North Holland Publishing Company.
\url{http://deleeuwpdx.net/janspubs/1977/chapters/deleeuw_C_77.pdf}.

\leavevmode\hypertarget{ref-deleeuw_A_84f}{}%
---------. 1984. ``Differentiability of Kruskals Stress at a Local
Minimum.'' \emph{Psychometrika} 49: 111--13.
\url{http://deleeuwpdx.net/janspubs/1984/articles/deleeuw_A_84f.pdf}.

\leavevmode\hypertarget{ref-hirriart-urruty_88}{}%
Hirriart-Urruty, J.-B. 1988. ``Generalized Differentiability / Duality
and Optimization for Problems Dealing with Differences of Convex
Functions.'' In \emph{Convexity and Duality in Optimization}, edited by
Ponstein. J., 37--70. Lecture Notes in Economics and Mathematical
Systems 256. Springer.

\leavevmode\hypertarget{ref-ivanov_16}{}%
Ivanov, V.I. 2016. ``Higher Order Optimality Conditions with an
Arbitrary Non-Differentiable Function.'' \emph{Optimization} 65 (11):
1909--27.

\leavevmode\hypertarget{ref-kruskal_64a}{}%
Kruskal, J.B. 1964. ``Multidimensional Scaling by Optimizing Goodness of
Fit to a Nonmetric Hypothesis.'' \emph{Psychometrika} 29: 1--27.

\leavevmode\hypertarget{ref-rockafellar_70}{}%
Rockafellar, R.T. 1970. \emph{Convex Analysis}. Princeton University
Press.

\leavevmode\hypertarget{ref-varga_62}{}%
Varga, R.S. 1962. \emph{Matrix Iterative Analysis}. Prentice-Hall.

\leavevmode\hypertarget{ref-vesely_zajicek_89}{}%
Vesely, L., and L. Zajicek. 1989. ``Delta-convex Mappings between Banach
Spaces and Applications.'' \emph{Dissertationes Mathematicae} 289:
1--48.


\end{document}
